%% ------------------------------------------------------------------------------------------------
%%  Abstract and keywords. Journal-agnostic.
%% ------------------------------------------------------------------------------------------------
\begin{abstract}
A group wants to choose one option by consensus. Anyone may propose options, every voter judges all of
them, and voters have a right of veto. We review what research says about voting in this setting and
derive a procedure that reaches the best compromise without ever getting stuck. The literature on
committees, veto players and consensus practice agrees that hard unanimity is rarely optimal: it breeds
delay, status-quo bias and holdouts, and in a large assembly it deadlocks almost surely, as we show for
spatial preferences. Social choice adds three constraints: no rule escapes manipulation, the majority
principle and a minority veto conflict in some profiles, and an open agenda requires a rule that is
independent of clones, which Borda's rule, fallback bargaining and the proportional veto core are not. We
propose \emph{Condorcet with a blocking quota}: on a single ballot, every voter ranks all options and marks
the unacceptable ones; options vetoed by at least a third of the voters are excluded, the least-vetoed are
kept if all are, and the Schulze winner among the rest is elected. The procedure always decides in one
round, gives the status quo no default advantage, elects an option acceptable to more than two thirds of
the voters whenever one exists, respects the majority among such options and ignores clones. In
simulations of polarised electorates it halves the share of contested outcomes relative to the Schulze
method, at no cost in total welfare. Its price is strategic blocking, which we quantify and which named,
reasoned vetoes and a bounded revision round help contain.

\medskip
\noindent\textbf{Keywords.} Voting rules, veto rights, consensus, Condorcet methods, Schulze method,
independence of clones, proportional veto, deadlock, compromise.
\end{abstract}
